\BourbakiExerciseGroup

\begin{enumerate}
\def\labelenumi{\arabic{enumi})}
\item
  Every topology compatible with the group structure of a finite group G may be obtained by taking as neighbourhoods of the identity element the sets containing a normal subgroup H of G.
\item
  A semi-topological group is a group \(\mathbf{G}\) endowed with a topology such that, for each \(a \in \mathbf{G}\) , the translations \(x \to ax\) and \(x \to xa\) are continuous on \(\mathbf{G}\) , and such that the symmetry \(x \to x^{-1}\) is continuous on \(\mathbf{G}(*)\) .
\end{enumerate}

\begin{enumerate}
\def\labelenumi{\alph{enumi})}
\tightlist
\item
  A filter \(\mathfrak{B}\) on a group \(\mathbf{G}\) is the neighbourhood filter of \(e\) in a topology which makes \(\mathbf{G}\) a semi-topological group if and only if \(\mathfrak{B}\) satisfies axioms \((\mathrm{GV}_{\mathrm{I}})\) and \((\mathrm{GV}_{\mathrm{III}})\) and the family of filters
\end{enumerate}

\[
\mathfrak {B} (x) = x. \mathfrak {B} = \mathfrak {B}. x \quad (x \in G)
\]

satisfies axiom \((\mathbf{V}_{\mathbf{IV}})\) of Chapter I, § 1, no. 2.

\begin{enumerate}
\def\labelenumi{\alph{enumi})}
\setcounter{enumi}{1}
\tightlist
\item
  If G is an infinite group, show that the topology in which the open sets are \(\varnothing\) and the complements of finite subsets of G makes G a non-Hausdorff semi-topological group in which \(\{e\}\) is closed, but is not compatible with the group structure of G.
\end{enumerate}

\(^{*}c)\) Define a topology on R, finer than the usual topology, for which R is a semi-topological but not topological group {[}consider a decreasing sequence \((r_{n})\) of numbers \textgreater0, tending to 0, and take the intersections of symmetric intervals {]}-a, a{[} and of the complement of the set of points \(\pm r_{n}\) ). *

(*) If G is locally compact these conditions imply that the topology of G is compatible with its group structure (Chapter X, § 3, Exercise 25).

\begin{enumerate}
\def\labelenumi{\arabic{enumi})}
\setcounter{enumi}{2}
\item
  Let A be a subset of a semi-topological (resp. topological) group G. If V runs through the neighbourhood filter of e in G, then the intersection of the sets AV, or of the sets VA (or of the sets VAV), is the closure \(\overline{A}\) of A.
\item
  A paratopological group is a group G endowed with a topology satisfying axiom \((\mathrm{GT}_{\mathbf{I}})(*)\) .
\end{enumerate}

\begin{enumerate}
\def\labelenumi{\alph{enumi})}
\item
  A filter \(\mathfrak{B}\) on a group \(G\) is the neighbourhood filter of \(e\) for a topology which makes \(G\) a paratopological group if and only if \(\mathfrak{B}\) satisfies axioms \((\mathrm{GV}_{\mathbf{I}})\) and \((\mathrm{GV}_{\mathbf{III}})\) and \(e \in V\) for all \(V \in \mathfrak{B}\) . The corresponding topology is then unique; it is Hausdorff if and only if the intersection of the sets \(V.V^{-1}\) , where \(V\) runs through \(\mathfrak{B}\) , consists of the point \(e\) alone.
\item
  On the group \(\mathbf{Z}\) of integers, let \(\mathbf{V}_n\) be the set consisting of \(0\) and all integers \(m \geqslant n\) , for each \(n > 0\) . Show that the \(\mathbf{V}_n\) form a fundamental system of neighbourhoods of \(0\) for a non-Hausdorff topology on \(\mathbf{Z}\) , and that with this topology \(\mathbf{Z}\) is a paratopological group in which \(\{0\}\) is closed.
\item
  A paratopological group G is a topological group if and only if, for each subset A of G, the intersection of the sets AV (or of the sets VA) is the closure \(\overline{A}\) of A (as V runs through the neighbourhood filter of e) (to see that the condition is sufficient, take A to be the complement of an open neighbourhood of e).
\end{enumerate}

\begin{enumerate}
\def\labelenumi{\arabic{enumi})}
\setcounter{enumi}{4}
\tightlist
\item
  Let \(\mathfrak{B}\) be a filter on a group \(\mathbf{G}\) , satisfying axioms \((\mathrm{GV}_{\mathbf{I}})\) and \((\mathrm{GV}_{\mathbf{II}})\) .
\end{enumerate}

\begin{enumerate}
\def\labelenumi{\alph{enumi})}
\tightlist
\item
  There is a unique topology \(\mathfrak{C}_s\) (resp. \(\mathfrak{C}_d\) ) on \(\mathbf{G}\) such that, for each \(a \in \mathbf{G}\) , the neighbourhood filter of \(a\) with respect to \(\mathfrak{C}_s\) (resp. \(\mathfrak{C}_d\) ) is \(a\) . \(\mathfrak{B}\) (resp. \(\mathfrak{B}\) . \(a\) ).
\end{enumerate}

Let \(\mathbf{G}_s\) (resp. \(\mathbf{G}_d\) ) be the topological space obtained by endowing \(\mathbf{G}\) with the topology \(\mathcal{C}_s\) (resp. \(\mathcal{C}_d\) ). For each \(a \in \mathbf{G}\) , the left translation \(x \to ax\) (resp. the right translation \(x \to xa\) ) is a continuous mapping of \(\mathbf{G}_s\) into \(\mathbf{G}_s\) (resp. of \(\mathbf{G}_d\) into \(\mathbf{G}_d\) ). The symmetry \(x \to x^{-1}\) is a continuous mapping of \(\mathbf{G}_s\) into \(\mathbf{G}_d\) and of \(\mathbf{G}_d\) into \(\mathbf{G}_s\) .

\begin{enumerate}
\def\labelenumi{\alph{enumi})}
\setcounter{enumi}{1}
\tightlist
\item
  The following conditions are equivalent: 1) V satisfies \((\mathrm{GV}_{\mathrm{III}})\) ;
\end{enumerate}

\begin{enumerate}
\def\labelenumi{\arabic{enumi})}
\setcounter{enumi}{1}
\tightlist
\item
  for each \(a \in G, x \to xa\) is a continuous mapping of \(G_{s}\) into \(G_{s}\) ;
\item
  \(x \to x^{-1}\) is a continuous mapping of \(G_{s}\) into \(G_{s}\) .
\end{enumerate}

\(^{*}c)\) Let \(G\) be the group \(\mathbf{GL}(2, \mathbf{Q}_p)\) of \(2 \times 2\) square matrices over the \(p\) -adic field \(\mathbf{Q}_p\) . For each integer \(n > 0\) , let \(H_n\) be the set of matrices of the form \(I + p^n U\) , where \(U\) is a \(2 \times 2\) matrix whose elements are \(p\) -adic integers. Show that the \(H_n\) are subgroups of \(G\) whose intersection consists only of the identity element. Deduce that if \(\mathfrak{B}\) is the filter which has the \(H_n\) as a base, then the corresponding topologies \(T_s\) and \(T_d\) on \(G\) are Hausdorff. Show that these topologies are distinct.*

\begin{enumerate}
\def\labelenumi{\arabic{enumi})}
\setcounter{enumi}{5}
\tightlist
\item
  Let G be a topological group and let V be an open neighbourhood of e in G. Show that the set of \(x \in G\) such that \(x^{2} \in V\) is the union of the neighbourhood W of e such that \(W^{2} \subset V\) .
\end{enumerate}

*7) Let \(\varphi\) be the canonical mapping of \(\mathbf{R}\) onto \(\mathbf{T} = \mathbf{R} / \mathbf{Z}\) , and let \(f\) be the mapping \(x \to \varphi(3x)\) , restricted to the neighbourhood \(\mathbf{V} = [-1/8, +1/8]\) of \(0\) in \(\mathbf{R}\) . Show that \(f\) is a homeomorphism of \(\mathbf{V}\) onto \(f(\mathbf{V})\) and satisfies condition 1) of Definition 2 of no. 3, but is not a local isomorphism from \(\mathbf{R}\) to \(\mathbf{T}\) . *

\begin{enumerate}
\def\labelenumi{\arabic{enumi})}
\setcounter{enumi}{7}
\item
  Let G be a Hausdorff topological group, and let \(s \neq e\) be a point of G. Show that there is a symmetric neighbourhood V of \(e\) such that \(s \notin V^2\) , and that G can be covered by at most 17 sets of the form \(a_k V b_k (1 \leqslant k \leqslant 17)\) . (Consider a maximal element in the set of symmetric neighbourhoods V of \(e\) such that \(s \notin V^2\) . We are thus led to consider the set S of \(x \in G\) such that \(x^2 = s\) ; observe that if \(x, y\) are two elements of S such that \(xy^{-1} \in S\) , then we must have \(x^2 = e\) , and if \(z\) is a third element of S such that \(xz^{-1} \in S\) and \(yz^{-1} \in S\) , then \(z\) must commute with \(xy^{-1}\) , and \(xy^{-1}z^{-1} \notin S\) .) If G is commutative, the number 17 may be replaced by 5.
\item
  Show that on a group G, the least upper bound of a family of topologies compatible with the group structure of G is compatible with the group structure of G.
\end{enumerate}

